\chapter{Semen Analysis}
\section*{What Is This Test?}
Semen analysis (seminogram) is the laboratory evaluation of semen to assess male fertility. It evaluates: semen volume, sperm concentration (count), total sperm count, sperm motility (percentage of progressively motile sperm), sperm morphology (percentage of normally shaped sperm by strict Kruger criteria), pH, viscosity, liquefaction time, and the presence of round cells (leukocytes vs. immature germ cells). Reference values are based on the World Health Organization (WHO) laboratory manual (5th Edition, 2010; 6th Edition, 2021), derived from a fertile population. Male factor infertility accounts for approximately 40--50\% of infertility cases. Post-vasectomy semen analysis is performed to confirm azoospermia after vasectomy. Semen is collected by masturbation after 2--7 days of sexual abstinence, directly into a sterile wide-mouth container, and delivered to the laboratory within 1 hour at body temperature (20--37 degrees C). Longer abstinence ($>$7 days) reduces motility; shorter (<2 days) reduces sperm count.
\section*{Alternative Names}
Semen Analysis; Seminogram; Sperm Count; Semen Profile; Ejaculate Analysis
\section*{Principle}
Semen analysis measures the cellular and noncellular components of the ejaculate. After liquefaction, the specimen volume, pH, viscosity, and liquefaction time are recorded, and the sample is examined microscopically: sperm concentration and total count are determined by manual counting in a hemocytometer or by computer-assisted semen analysis (CASA), motility is graded as progressive, non-progressive, or immotile, and morphology is scored on a stained smear using strict Kruger criteria. Post-vasectomy specimens are centrifuged and the pellet is examined in its entirety to confirm azoospermia.
\section*{History}
Spermatozoa were first observed by Antonie van Leeuwenhoek, with his assistant Johan Ham, in 1677, who described them in human semen. Systematic quantitative evaluation of semen emerged in the 1920s--1930s as the specialty of infertility medicine developed. The World Health Organization published its first laboratory manual for the examination of human semen in 1980, establishing standardized methods so that results could be compared across laboratories. The WHO 5th edition (2010) redefined the reference ranges using data from approximately 1,900 fertile men across five continents, replacing fixed ``normal'' values with 5th centile lower reference limits. The 6th edition (2021) updated these limits and incorporated computer-assisted semen analysis and newer techniques for assessing sperm DNA integrity. The WHO manual remains the international standard for performing and interpreting semen analysis.
\section*{Why This Test Is Ordered}
\begin{itemize}
  \item Evaluation of male infertility: the initial investigation in couples unable to conceive after 12 months of regular unprotected intercourse (or 6 months when the female partner is 35 years or older)
  \item Confirmation of successful vasectomy: demonstration of azoospermia on post-vasectomy analysis at 8--16 weeks
  \item Assessment after testicular trauma, orchitis, torsion, or exposure to gonadotoxic agents (chemotherapy, radiation)
  \item Evaluation of hypogonadism or suspected endocrine causes of infertility, in conjunction with hormonal testing
  \item Diagnosis of retrograde ejaculation or anejaculation (aspermia), including in men with spinal cord injury or diabetic neuropathy
  \item Baseline assessment before assisted reproductive technology (IVF, ICSI) or sperm banking
  \item Investigation of ejaculatory disorders, including hematospermia
\end{itemize}
\section*{Primary vs Secondary Test}
Semen analysis is the primary diagnostic test for male factor infertility and the central first-line investigation in the evaluation of the infertile couple. It is also the definitive confirmatory test after vasectomy. Because of substantial intra-individual variability, a single abnormal result is not diagnostic on its own and should be confirmed by repeat testing.
\section*{How This Test Is Performed}
Semen is collected by masturbation after 2--7 days of sexual abstinence into a sterile wide-mouth container and delivered to the laboratory within 1 hour, kept at body temperature (20--37 degrees C). In the laboratory, the specimen is allowed to liquefy (normally within 30 minutes) and is assessed for volume, pH, viscosity, and liquefaction time; sperm concentration, total count, and motility are determined microscopically or by computer-assisted semen analysis (CASA), and morphology is scored on stained smears using strict Kruger criteria. Results are typically available within 1--2 days.
\section*{What Preparation Is Needed}
Sexual abstinence of 2--7 days before collection is required; abstinence shorter than 2 days lowers sperm count, and longer than 7 days lowers motility. No fasting is needed. The specimen must be collected directly into the sterile container supplied by the laboratory, avoiding lubricants (which may be spermicidal), and delivered within 1 hour. Recent febrile illness, acute infection, and certain medications (e.g., testosterone, sulfasalazine) can transiently depress sperm parameters.
\section*{Contraindications and Precautions}
\begin{itemize}
  \item There are no absolute contraindications. Precautions: the ejaculate must be complete, since loss of the sperm-rich initial portion invalidates the analysis; specimens collected during intercourse require special nonspermicidal collection condoms; extremes of temperature or delayed transport degrade motility; and a single abnormal result should be confirmed with a repeat analysis after 2--3 months, as sperm parameters fluctuate significantly between samples.
\end{itemize}
\section*{Risks}
No physical risks are associated with the test itself. Some men find collection by masturbation embarrassing, and alternative collection methods are available. No venipuncture is involved.
\section*{What Happens After the Test}
Results are usually available within 1--2 days and are compared against the WHO lower reference limits. Abnormal parameters are confirmed by repeat testing after 2--3 months and prompt further evaluation (hormonal studies, scrotal ultrasound, genetic testing). After vasectomy, azoospermia on post-vasectomy samples confirms successful sterilization, but patients should continue contraceptive precautions until this is documented.
\section*{Understanding Results}
Oligozoospermia: Sperm concentration <15 million/mL. Asthenozoospermia: Progressive motility <32\%. Teratozoospermia: Normal morphology <4\%. Oligoasthenoteratozoospermia (OAT): All three parameters reduced (most common pattern). Azoospermia: No sperm in the ejaculate after centrifugation --- obstructive (normal FSH/testosterone, congenital bilateral absence of vas deferens [CBAVD], vasectomy) vs. non-obstructive (elevated FSH, reduced testosterone, Sertoli cell only syndrome, Klinefelter syndrome 47,XXY). Aspermia: No ejaculate (anejaculation) --- retrograde ejaculation (sperm found in post-ejaculatory urine), anejaculation (spinal cord injury, diabetic neuropathy, medications [alpha-blockers, antipsychotics]).
\section*{Normal Results}
WHO 6th Edition (2021) lower reference limits (5th centile): Volume $\geq$1.4 mL. Sperm concentration $\geq$16 million/mL. Total sperm count $\geq$39 million/ejaculate. Progressive motility $\geq$30\%. Total motility (progressive + non-progressive) $\geq$42\%. Normal morphology $\geq$4\% (strict Kruger criteria). Vitality (live sperm) $\geq$54\%. White blood cells <1.0 million/mL. pH 7.2--8.0. Liquefaction within 30 minutes.
\section*{Follow-up Tests}
Repeat semen analysis after 2--3 months (a full spermatogenic cycle is ~74 days). Hormonal evaluation: serum FSH, LH, testosterone, prolactin. Scrotal ultrasound. Genetic testing: karyotype, Y-chromosome microdeletion analysis, CFTR gene testing (CBAVD). Post-ejaculatory urinalysis (retrograde ejaculation). Sperm DNA fragmentation assay. Anti-sperm antibody testing. Post-vasectomy semen analysis at 8--16 weeks post-procedure.
\section*{References}
\begin{enumerate}
  \item World Health Organization. WHO Laboratory Manual for the Examination and Processing of Human Semen. 6th ed. WHO Press; 2021.
  \item Cooper TG, Noonan E, von Eckardstein S, et al. World Health Organization reference values for human semen characteristics. Hum Reprod Update. 2010;16(3):231--245.
  \item World Health Organization. WHO Laboratory Manual for the Examination and Processing of Human Semen. 5th ed. WHO Press; 2010.
  \item Practice Committee of the American Society for Reproductive Medicine. Diagnostic evaluation of the infertile male: a committee opinion. Fertil Steril. 2015;103(3):e18--e25.
  \item Esteves SC, Zini A, Aziz N, et al. Critical appraisal of World Health Organization's new reference values for human semen characteristics and effect on diagnosis and treatment of subfertile men. Urology. 2012;79(1):16--22.
\end{enumerate}
\clearpage
\section*{Regulatory Status and Authorization Date}
This chapter describes a test category, not a specific commercial product. FDA, CDSCO, EU IVDR, and MHRA approval or registration dates therefore cannot be assigned without the assay manufacturer, product name, instrument, and intended use. For laboratory-developed tests, an individual agency approval date may not exist. See the Preface for the interpretation of regulatory status.

\clearpage
